\documentclass[lettersize,journal]{IEEEtran}
\usepackage{amsmath,amsfonts}
\usepackage{algorithmic}
\usepackage{algorithm}
\usepackage{array}
\usepackage[caption=false,font=normalsize,labelfont=sf,textfont=sf]{subfig}
\usepackage{textcomp}
\usepackage{stfloats}
\usepackage{url}
\usepackage{verbatim}
\usepackage{graphicx}
\usepackage{cite}
\usepackage{booktabs}
\usepackage{multirow}
\usepackage{amssymb}
\hyphenation{op-tical net-works semi-conduc-tor IEEE-Xplore}

\begin{document}

\title{Structure Information Graph Learning with Graph Convolutional Networks\\ for Anomaly Detection in Mixed-Attribute Data}

\author{Ziwei~Cui\textsuperscript{1},~Jiaojiao~Niu\textsuperscript{1*},~
        Degang~Chen\textsuperscript{2},~and~Jinhai~Li\textsuperscript{3,4}
\thanks{\textsuperscript{1}School of Computer Science, Yangtze University, Jingzhou, Hubei 434023, China.}
\thanks{\textsuperscript{2}School of Mathematics and Physics, North China Electric Power University, Beijing 102206, China.}
\thanks{\textsuperscript{3}Yunnan Key Laboratory of Complex Systems and Brain-Inspired Intelligence,
        Kunming University of Science and Technology, Kunming 650500, China.}
\thanks{\textsuperscript{4}Faculty of Science, Kunming University of Science and Technology, Kunming 650500, China.}
\thanks{*Corresponding author: Jiaojiao Niu.}
\thanks{Manuscript received xxx; revised xxx.}}

\markboth{IEEE Transactions on Knowledge and Data Engineering,~Vol.~XX, No.~XX,~2026}%
{Cui \MakeLowercase{\textit{et al.}}: NMIGOD for Anomaly Detection in Mixed-Attribute Data}

\maketitle

\begin{abstract}
Anomaly detection in mixed-attribute data remains a challenging problem due to heterogeneous feature types
and the difficulty of modeling structural relationships among samples. Traditional distance-based and
neighborhood-based methods often fail to capture complex dependencies between samples, leading to
unsatisfactory detection performance. To address this issue, this paper proposes a structure information
graph learning framework for anomaly detection in mixed data, termed NMIGOD. Specifically, a Neighborhood
Mutual Information (NMI)-based measure is introduced to characterize structural dependencies by modeling
the consistency of local neighborhood structures. Based on this measure, a weighted graph is constructed
where edge weights encode neighborhood-level structural similarity rather than point-wise distance. A
threshold-based edge filtering strategy is then applied to remove weak structural dependencies to generate
a compact structure-preserving graph. On top of the constructed graph, a Graph Convolutional Network is
employed to propagate structural information and learn discriminative representations for anomaly scoring
under weakly supervised settings. Extensive experiments on 24 UCI benchmark datasets demonstrate that
the proposed method consistently outperforms state-of-the-art approaches. The code is available at:
\url{https://github.com/zhiyuandaoren/anomaly-detection-NMIGOD.git}.
\end{abstract}

\begin{IEEEkeywords}
Anomaly detection, neighborhood mutual information, graph structure learning, graph convolutional networks, mixed-attribute data.
\end{IEEEkeywords}

% ===================================================================
\section{Introduction}
% ===================================================================
\IEEEPARstart{A}{nomaly} detection aims to identify rare patterns or observations that deviate markedly
from the majority of data \cite{ref1,ref2}. Since anomalies often signal critical events, it has become
a fundamental task in many high-stakes domains, such as fraud detection and industrial fault diagnosis
\cite{ref3,ref4}. However, real-world data often exhibit high complexity, particularly in the form of
mixed-attribute data, making it difficult for traditional distance-based or density-based methods to
provide a unified similarity measure. As structural information offers a more robust basis for similarity
assessment by capturing how objects are situated within their local neighborhoods, it is meaningful to
design anomaly detection methods based on structural similarity rather than conventional feature distances.

Existing anomaly detection methods can be broadly divided into three categories. The first category
comprises traditional approaches, including statistical models \cite{ref5}, distance-based methods such
as $k$-nearest neighbors ($k$-NN) \cite{ref6} and local outlier factor (LOF) \cite{ref7}, and
clustering-based techniques \cite{ref8}. These methods identify anomalies by measuring deviations in the
original feature space through probabilistic assumptions, geometric distance, or cluster structures.
Although these methods are simple and interpretable, the effectiveness is often limited in complex
real-world scenarios, particularly when data exhibits mixed types or high-dimensional structures. The
second category consists of methods that incorporate other theoretical tools, such as fuzzy sets, rough
sets \cite{ref11} and formal concept analysis, to better handle data uncertainty and heterogeneity. Among
these, neighborhood rough sets (NRS) \cite{ref12,ref13} have gained particular attention due to their
data-driven nature and the ability to handle mixed attributes without requiring prior probability
distributions. Representative NRS-based works include NIEOD \cite{ref14}, which defines neighborhood
information entropy and uses its variation for outlier scoring; ADFNR \cite{ref16} introduces fuzzy
relations into the NRS framework to handle noise and uncertainty; GCOD \cite{ref19} is the first to
apply formal concept analysis to outlier detection by measuring distances between granular concepts, and
DASOD \cite{ref20} detects outliers from both object and attribute perspectives within the FCA framework
and integrates structural deviation with attribute-level rarity. These methods are effective at
characterizing local uncertainty, but they operate on individual granules and fail to capture global
topological dependencies among data objects. The third category encompasses deep learning-based methods,
particularly graph neural networks (GNNs) and graph convolutional networks (GCNs) \cite{ref21}, which
have been widely adopted for their ability to learn from relational structures. Notable examples include
LUNAR \cite{ref22}, DOMINANT \cite{ref23}, and ADA-GAD \cite{ref24}. However, these methods typically
construct graphs using conventional feature-space distance (e.g., Euclidean distance or $k$-NN rules),
which are unreliable for mixed-attribute data. Moreover, under the multi-layer message passing mechanism
of GCNs, outliers that are erroneously connected to normal clusters tend to have their anomaly signals
averaged out by numerous normal neighbors, leading to the well-known over-smoothing problem \cite{ref25}.
As a result, NRS-based methods capture local structures but lack global context, while GCN-based approaches
model global dependencies but rely on distance-based graph construction, which may introduce spurious
connections and dilute anomaly signals in mixed-attribute settings.

As stated above, both NRS-based and GCN-based approaches have complementary strengths and weaknesses,
making their integration a natural strategy. For achieving this, the key issue is how to construct a graph
that captures local neighborhood information while supporting global learning. Existing graph construction
methods can be broadly categorized into heuristic-based and learning-based
approaches~\cite{Mlyahilu2025}, with the former typically relying on feature-space distances (e.g.,
Euclidean or $k$-NN) or information-theoretic measures such as transfer
entropy~\cite{Febrinanto2025}, and the latter learning graph topologies adaptively from
data~\cite{He2024,Pang2024}. However, these approaches still fundamentally operate in the feature space,
without explicitly considering structural similarity between local neighborhoods. To address this, this
paper proposes a structure information graph learning approach that shifts the basis of graph construction
from feature distance to structural similarity, where two objects are connected if they share similar
local neighborhood configurations. To quantify such structural similarity, Neighborhood Mutual Information
(NMI) within the neighborhood rough set framework is introduced. A critical issue in graph learning is
the determination of neighborhood scales, which significantly affects the quality of the resulting graph.
In the proposed model, an entropy-based adaptive radius is introduced, eliminating the need for manual
parameter tuning. The constructed graph is then fed into a GCN for semi-supervised anomaly detection.
Extensive experiments on 24 UCI benchmark datasets, covering various data types and anomaly ratios, are
conducted to evaluate the proposed method against several state-of-the-art baselines. The experimental
results demonstrate that the proposed method consistently achieves superior performance across multiple
evaluation metrics. The main contributions of this paper are summarized as follows:

\begin{enumerate}
\item An entropy-based adaptive radius is proposed to determine neighborhood scales without manual tuning;
\item A structure information graph is constructed using Neighborhood Mutual Information to capture
      structural similarity;
\item A novel anomaly detection framework named NMIGOD is proposed by integrating NMI-based graph
      construction with GCN for semi-supervised learning.
\end{enumerate}

The remainder of this paper is organized as follows: Section~2 reviews the preliminaries. Section~3
describes the proposed NMIGOD framework in detail. Section~4 reports the experimental results, and
Section~5 concludes the paper.

% ===================================================================
\section{Preliminaries}
% ===================================================================
This section reviews two essential preliminaries: neighborhood information systems and graph
convolutional networks.

\subsection{Neighborhood Information System}

A neighborhood information system is an extension of the classical information system by incorporating
neighborhood relations. Therefore, the classical information system is first recalled as follows.

An information system is formally defined as a quadruple \(IS = (U, A, V, f)\), where
\(U = \{x_{1}, x_{2}, \cdots, x_{n}\}\) is a non-empty finite set of objects,
\(A = \{a_{1}, a_{2}, \cdots, a_{m}\}\) is a non-empty finite set of attributes,
\(V = \bigcup_{a \in A} V_{a}\) is the union of attribute domains \(V_{a}\), and
\(f: U \times A \rightarrow V\) denotes an information function such that \(f(x,a) \in V_{a}\) for all
\(x \in U\) and \(a \in A\).

The construction of neighborhoods often relies on a distance function defined over the attribute set,
which serves as the basis for granulating the universe and establishing neighborhood relations.

\textbf{Definition 1} \cite{ref14}. Let \(IS = (U, A, V, f)\) be a neighborhood information system.
For \(x \in U\), an attribute subset \(B \subseteq A\), and a radius \(\varepsilon \geq 0\), the
\(\varepsilon\)-neighborhood of object \(x\) with respect to \(B\) is defined as:
\begin{equation}
N_{B}^{\varepsilon}(x) = \left\{ y \in U \mid d_{B}(x, y) \leq \varepsilon \right\},
\end{equation}
where \(d_{B}: U \times U \rightarrow \mathbb{R}^{+}\) is a distance function. That is,
\(N_{B}^{\varepsilon}(x)\) collects all objects whose distance from \(x\) on the attribute subset \(B\)
does not exceed \(\varepsilon\).

The neighborhood relation induced on \(U\) by \(B\) is defined as:
\begin{equation}
NR_{B}^{\varepsilon} = \{(x, y) \in U \times U \mid y \in N_{B}^{\varepsilon}(x)\}.
\end{equation}

This relation is reflexive and symmetric. A neighborhood cover (also known as neighborhood knowledge)
of \(U\) is constituted by the quotient set
\(U / NR_{B}^{\varepsilon} = \{N_{B}^{\varepsilon}(x) \mid x \in U\}\), and the underlying granular
structure for subsequent analysis is established based on this cover.

A neighborhood information system (NIS) is then defined as
\(NIS = (U, NR_{A}^{\varepsilon}, V, f)\), where \(NR_{A}^{\varepsilon}\) denotes the neighborhood
relation induced by the full attribute set \(A\). When a decision attribute \(D\) is present, the system
can extend to a neighborhood decision system (NDS):
\(NDS = (U, NR_{C}^{\varepsilon} \cup D, V, f)\), where \(C \subseteq A\) denotes the set of conditional
attributes and \(D\) is the decision attribute.

\subsection{Graph Convolutional Networks}

Graph Convolutional Networks (GCNs) \cite{ref21} extend the convolution operation from image data to
graph-structured data. Given a graph with adjacency matrix \(M \in \mathbb{R}^{n \times n}\) and node
feature matrix \(X \in \mathbb{R}^{n \times m}\), each GCN layer updates node representations by
aggregating information from neighboring nodes:
\begin{equation}
H^{(l+1)} = \sigma\left( \widehat{M} H^{(l)} W^{(l)} \right),
\end{equation}
where \(\widehat{M} = \widetilde{D}^{-\frac{1}{2}} \widetilde{M} \widetilde{D}^{-\frac{1}{2}}\) denotes
the symmetrically normalized adjacency matrix with added self-loops, i.e.,
\(\widetilde{M} = M + I_{n}\), and \(\widetilde{D}\) is the corresponding degree matrix.
\(H^{(l)}\) are the node representations at layer \(l\), with \(H^{(0)} = X\), and \(W^{(l)}\) is a
trainable weight matrix. \(\sigma\) is a non-linear activation function, typically set to ReLU.

For a two-layer GCN, the forward propagation is:
\begin{equation}
Z = \text{softmax}\left( \widehat{M} \cdot \text{ReLU}\left( \widehat{M} X W^{(0)} \right) W^{(1)} \right),
\end{equation}
where \(Z \in \mathbb{R}^{n \times 2}\) produces a probability distribution over the normal and anomalous classes.

In the semi-supervised setting, the model is trained with a masked cross-entropy loss that only accounts
for labeled nodes:
\begin{equation}
\mathcal{L} = \sum_{i \in \mathcal{L}_{L}} \sum_{c=1}^{2} Y_{ic} \ln Z_{ic},
\end{equation}
where \(\mathcal{L}_{L}\) denotes the set of labeled nodes, \(Y_{ic}\) is the ground-truth label indicator
(i.e., \(Y_{ic} = 1\) if node \(i\) belongs to class \(c\), and otherwise 0), and \(Z_{ic}\) is the
predicted probability that node \(i\) belongs to class \(c\).

% ===================================================================
\section{Proposed Method}
% ===================================================================
Building upon the neighborhood information system and GCN introduced in Section~2, this section presents
the proposed anomaly detection framework. The core idea is to construct a graph based on structural
similarity---quantified by Neighborhood Mutual Information (NMI)---rather than conventional feature
distance, and then perform semi-supervised anomaly detection on this graph using a two-layer GCN.
As illustrated in Fig.~\ref{fig:framework}, the framework consists of three stages: (1) neighborhood
system construction; (2) NMI-based graph construction; and (3) GCN-based anomaly detection. The
remainder of this section details each stage in sequence, followed by the algorithm summary and
complexity analysis.

\begin{figure}[!t]
\centering
\includegraphics[width=\columnwidth]{framework.pdf}
\caption{Overview of the proposed framework.}
\label{fig:framework}
\end{figure}

\subsection{Neighborhood System Construction}

Given an information system \(IS = (U, A, V, f)\), the construction of neighborhood knowledge requires
a distance metric and a set of per-attribute thresholds. In Section~2.1, the classical neighborhood
system is defined based on a global distance function \(d_{A}\) over the full attribute set \(A\) with a
unified radius \(\varepsilon\). However, when dealing with mixed-attribute data containing both numerical
and categorical features, such a global formulation may fail to capture the heterogeneous characteristics
of different attributes. To address this limitation, the Heterogeneous Euclidean-Overlap Metric (HEOM)
\cite{ref15} is adopted as the distance function.

For \(x, y \in U\), the HEOM distance \(\text{HEOM}_{A}(x, y)\) with respect to \(A\) is defined as:
\begin{equation}
\text{HEOM}_{A}(x, y) = \sqrt{\sum_{a \in A} d_{a}(x, y)^{2}},
\end{equation}
where \(d_{a}(x, y)\) denotes the component-wise distance under attribute \(a\), given by:
\begin{equation}
d_{a}(x, y) = \begin{cases}
0, & \text{if } a \text{ is categorical and } f(x,a) = f(y,a); \\
1, & \text{if } a \text{ is categorical and } f(x,a) \neq f(y,a); \\
|f(x,a) - f(y,a)|, & \text{if } a \text{ is numerical}.
\end{cases}
\end{equation}

The distance \(d_{a}(x, y)\) is bounded in \([0, 1]\), where 0 indicates that two objects are identical
or highly similar on attribute \(a\), and 1 indicates they are completely different or incomparable. To
ensure this, all numerical attributes are preprocessed using min-max normalization before distance
computation.

Instead of relying on a single global radius for all attributes, we compute an adaptive radius for each
numerical attribute. The radius is determined based on the standard deviation of the attribute and an
entropy-based adjustment factor.

\textbf{Definition 2.} For a numerical attribute \(a \in A\), let the objects in \(U\) be sorted in
ascending order of their values on \(a\): \(v_{1} \leq v_{2} \leq \cdots \leq v_{n}\), where
\(v_{i} = f(x_{i}, a)\). The gaps between adjacent values are computed as
\(\text{gap}_{i} = v_{i+1} - v_{i},\; i = 1, 2, \cdots, n-1\). Using the standard deviation
\(\text{std}(a)\) as the threshold, the sorted sequence is split into connected components
\(C = \{C_{1}, C_{2}, \cdots, C_{k}\}\) by cutting at positions where \(\text{gap}_{i} > \text{std}(a)\).
The attribute neighborhood entropy is then defined as:
\begin{equation}
NE(a) = -\sum_{C_{i} \in C} \frac{|C_{i}|}{|U|} \log_{2} \frac{|C_{i}|}{|U|}.
\end{equation}

This entropy measures the uniformity of the attribute's value distribution after partitioning by the
standard deviation-based gap threshold. Based on this entropy, the density factor \(\rho_{a}\) is
computed as:
\begin{equation}
\rho_{a} = 1 - \frac{NE(a)}{\log_{2}|U|},
\end{equation}
which ranges from 0 to 1. The adaptive radius \(\varepsilon_{a}\) for attribute \(a\) is then defined as:
\begin{equation}
\varepsilon_{a} = \begin{cases}
0, & \text{if } a \text{ is categorical}; \\[6pt]
\displaystyle\frac{\text{std}(a)}{1 + \rho_{a}}, & \text{if } a \text{ is numerical}.
\end{cases}
\end{equation}
where \(\text{std}(a)\) represents the standard deviation of attribute \(a\). The overall procedure is
summarized in Algorithm~1.

\begin{algorithm}[!t]
\caption{Neighborhood and Adaptive Radius Construction (NARC)}
\label{alg:narc}
\begin{algorithmic}
\STATE \textbf{Input:} Information system \(IS = (U, A, V, f)\), where \(A = A_{num} \cup A_{cat}\),
       with \(A_{num}\) and \(A_{cat}\) being the sets of numerical and categorical attributes,
       respectively, \(n = |U|\).
\STATE \textbf{Output:} Adaptive radius set \(\varepsilon = \{\varepsilon_{a} \mid a \in A\}\),
       neighborhood knowledge \(N = \{N_{a}(x) \mid x \in U, a \in A\}\).
\STATE
\STATE // Min-max normalization for numerical attributes
\FOR{each \(a \in A_{num}\)}
    \STATE \(\max_{a} \leftarrow \max\{f(x,a) \mid x \in U\}\)
    \STATE \(\min_{a} \leftarrow \min\{f(x,a) \mid x \in U\}\)
    \FOR{each \(x \in U\)}
        \STATE \(f(x,a) \leftarrow (f(x,a) - \min_{a}) / (\max_{a} - \min_{a})\)
    \ENDFOR
\ENDFOR
\STATE
\STATE // Compute standard deviation for each numerical attribute
\FOR{each \(a \in A_{num}\)}
    \STATE \(\sigma_{a} \leftarrow \text{std}(f(x,a) \mid x \in U)\)
    \IF{\(\sigma_{a} = 0\)} \STATE \(\sigma_{a} \leftarrow 10^{-6}\) \ENDIF
\ENDFOR
\STATE
\STATE // Compute attribute neighborhood entropy and adaptive radius
\FOR{each \(a \in A_{num}\)}
    \STATE Sort \(U\) by \(f(x,a)\) ascending: \(v_{1} \leq v_{2} \leq \cdots \leq v_{n}\)
    \STATE Compute gaps: \(\text{gap}_{i} \leftarrow v_{i+1} - v_{i}\) for \(i = 1, \dots, n-1\)
    \STATE Partition sorted \(U\) into components \(C = \{C_{1}, \dots, C_{k}\}\) where consecutive
            objects belong to the same component if their gap does not exceed \(\sigma_{a}\)
    \STATE Compute \(NE(a)\) by Eq.~(8)
    \STATE Compute \(\varepsilon_{a}\) by Eq.~(10)
\ENDFOR
\STATE
\STATE // Build final neighborhoods
\FOR{each \(x \in U\)}
    \FOR{each \(a \in A\)}
        \IF{\(a \in A_{num}\)}
            \STATE \(N_{a}(x) = \{y \in U \mid |f(x,a) - f(y,a)| \leq \varepsilon_{a}\}\)
        \ELSE
            \STATE \(N_{a}(x) = \{y \in U \mid f(x,a) = f(y,a)\}\)
        \ENDIF
    \ENDFOR
\ENDFOR
\STATE \textbf{return} \(\varepsilon, N\)
\end{algorithmic}
\end{algorithm}

\subsection{NMI-Based Graph Construction}

In anomaly detection, whether an object is abnormal depends not only on its individual attribute values,
but also on its relational context with other objects. An object that deviates from the local structural
patterns of the majority is more likely to be an outlier. Therefore, it is necessary to characterize the
structural relationships between objects in a way that captures how similarly they are embedded in the
local neighborhood structure. For achieving this, a graph is constructed by considering the structural
similarity between objects. Before constructing the graph, however, a measure is required to quantify
such similarity between two objects.

A natural intuition is that two objects should be considered structurally similar if their neighborhoods
share common elements. This idea is analogous to the concept of mutual information in information theory,
where two random variables are considered correlated if they share information. Motivated by this, the
Neighborhood Mutual Information (NMI) between two objects is defined to measure the shared structural
information between their neighborhood knowledge. To formally define NMI, the neighborhood information
entropy is first introduced as follows.

\textbf{Definition 3.} Given the neighborhood knowledge
\(\mathbb{N}(x) = \{N^{\varepsilon_{a}}(x) \mid a \in A\} = \{N_{1}, \ldots, N_{k}\}\) of object
\(x \in U\) with respect to \(A\), the object neighborhood entropy of \(x\) is defined as:
\begin{equation}
NE(x) = -\sum_{i \in [1, k]} \frac{|N_{i}|}{|U|} \log_{2} \frac{|N_{i}|}{|U|}.
\end{equation}

\textbf{Definition 4.} For objects \(x, y \in U\), the neighborhood joint entropy with respect to
attribute set \(A\) is defined as:
\begin{equation}
NJE(x, y) = -\sum_{N_{i} \in \mathbb{N}(x)} \sum_{N_{j} \in \mathbb{N}(y)}
            \frac{|N_{i} \cap N_{j}|}{|U|} \log_{2} \frac{|N_{i} \cap N_{j}|}{|U|}.
\end{equation}

\textbf{Definition 5.} The neighborhood mutual information between two objects \(x, y \in U\) with
respect to attribute set \(A\) is defined as:
\begin{equation}
NMI(x, y) = NE(x) + NE(y) - NJE(x, y).
\end{equation}

The expanded form is:
\begin{equation}
NMI(x, y) = -\sum_{N_{i} \in \mathbb{N}(x)} \sum_{N_{j} \in \mathbb{N}(y)}
            \frac{|N_{i} \cap N_{j}|}{|U|}
            \log_{2} \frac{|N_{i} \cap N_{j}| \cdot |U|}{|N_{i}| \cdot |N_{j}|},
\end{equation}
with the convention that \(0 \log_{2} 0 = 0\).

The NMI measure satisfies the following properties:

\noindent \textbf{(1) Non-negativity.} \(NMI(x, y) \geq 0\), with equality if and only if
\(N_{i} \cap N_{j} = \varnothing\) for all \(N_{i} \in \mathbb{N}(x)\) and
\(N_{j} \in \mathbb{N}(y)\), i.e., when the neighborhood structures of \(x\) and \(y\) have no overlap
across any attribute.

\noindent \textbf{(2) Symmetry.} \(NMI(x, y) = NMI(y, x)\), which follows directly from the symmetric
construction and ensures that the resulting graph is undirected.

\noindent \textbf{(3) Structural discriminability.} A larger NMI value indicates stronger structural
similarity between the two objects, suggesting that they reside in comparable local contexts across
multiple attribute views. In contrast, a small NMI implies that the two objects have distinct
neighborhood configurations, which is typical for anomalous instances that deviate from the majority.

Based on NMI, a weighted undirected graph \(G = (V, E, W)\) can be constructed, where each object
\(x \in U\) serves as a node. The weight \(w_{xy}\) of the edge between two nodes \(x\) and \(y\) is
defined as:
\begin{equation}
w_{xy} = \frac{NMI(x, y)}{\sqrt{NE(x) \cdot NE(y)}}.
\end{equation}

Obviously, the edge weight is confined to the unit interval. This follows from the fact that
\(NMI(x, y) \leq \min\{NE(x), NE(y)\} \leq \sqrt{NE(x) \cdot NE(y)}\), it follows that
\(w_{xy} \in [0, 1]\), and specifically \(w_{xx} = 1\) when \(x = y\).

To reduce graph density and suppress noisy connections, a threshold-based sparsification strategy is
applied. The entries of the resulting sparse adjacency matrix \(M\):
\begin{equation}
M_{xy} = \begin{cases}
w_{xy}, & \text{if } w_{xy} \geq \tau, \\
0, & \text{otherwise}.
\end{cases}
\end{equation}
where \(\tau \in [0, 1]\) is a pre-defined threshold. The NMI-based graph construction process is
detailed in Algorithm~2.

\begin{algorithm}[!t]
\caption{NMI-based Graph Construction (NGC)}
\label{alg:ngc}
\begin{algorithmic}
\STATE \textbf{Input:} Information system \(IS = (U, A, V, f)\), \(n = |U|\),
       \(\varepsilon, N\) (from Algorithm~1), threshold \(\tau\)
\STATE \textbf{Output:} Normalized adjacency matrix \(M \in \mathbb{R}^{n \times n}\)
\STATE
\STATE // Compute neighborhood information entropy \(NE(x)\) (Eq.~11)
\FOR{each \(x \in U\)}
    \STATE \(NE(x) \leftarrow 0\)
    \FOR{each \(a \in A\)}
        \STATE \(p \leftarrow |N_{a}(x)| / n\)
        \IF{\(p > 0\)}
            \STATE \(NE(x) \leftarrow NE(x) - p \cdot \log_{2}(p)\)
        \ENDIF
    \ENDFOR
\ENDFOR
\STATE
\STATE // Compute Neighborhood Mutual Information \(NMI(x, y)\) (Eq.~14)
\FOR{each \(x_{i}, x_{j} \in U\)}
    \STATE \(NMI(x_{i}, x_{j}) \leftarrow 0\)
    \FOR{each \(a \in A\)}
        \STATE \(I_{a} \leftarrow |N_{a}(x_{i}) \cap N_{a}(x_{j})|\)
        \IF{\(I_{a} > 0\)}
            \STATE \(\text{ratio} \leftarrow (I_{a} \cdot n) / (|N_{a}(x_{i})| \cdot |N_{a}(x_{j})|)\)
            \STATE \(NMI(x_{i}, x_{j}) \leftarrow NMI(x_{i}, x_{j}) + (I_{a} / n) \cdot \log_{2}(\text{ratio})\)
        \ENDIF
    \ENDFOR
\ENDFOR
\STATE // Set self-NMI = object entropy
\FOR{each \(x_{i} \in U\)}
    \STATE \(NMI(x_{i}, x_{i}) \leftarrow NE(x_{i})\)
\ENDFOR
\STATE
\STATE // Compute edge weights (Eq.~15)
\FOR{each \(x_{i}, x_{j} \in U\)}
    \STATE \(w_{ij} \leftarrow NMI(x_{i}, x_{j}) / \sqrt{NE(x_{i}) \cdot NE(x_{j})}\)
\ENDFOR
\STATE // Set self-loops to 1
\FOR{each \(x_{i} \in U\)} \STATE \(w_{ii} \leftarrow 1\) \ENDFOR
\STATE
\STATE // Threshold sparsification (Eq.~16)
\FOR{each \(i, j \in 1, \dots, n\)}
    \IF{\(w_{ij} < \tau\)} \STATE \(M_{ij} \leftarrow 0\)
    \ELSE \STATE \(M_{ij} \leftarrow w_{ij}\) \ENDIF
\ENDFOR
\STATE
\STATE // Symmetric normalization \(D^{-1/2} M D^{-1/2}\)
\FOR{each \(i \in 1, \dots, n\)}
    \STATE \(d_{i} \leftarrow \sum_{j=1}^{n} M_{ij}\)
    \IF{\(d_{i} > 0\)} \STATE \(D_{\text{inv}, ii} \leftarrow d_{i}^{-1/2}\)
    \ELSE \STATE \(D_{\text{inv}, ii} \leftarrow 0\) \ENDIF
\ENDFOR
\STATE \(M \leftarrow D_{\text{inv}} \cdot M \cdot D_{\text{inv}}\)
\STATE \textbf{return} \(M\)
\end{algorithmic}
\end{algorithm}

\subsection{GCN-Based Anomaly Detection}

With the sparse graph \(M \in \mathbb{R}^{n \times n}\) constructed in the previous stage, anomaly
detection is performed using a two-layer GCN in a semi-supervised manner. The GCN propagation rule and
masked cross-entropy loss introduced in Section~2.2 are applied, with the adjacency matrix set to
\(\widehat{M} = \widetilde{D}^{-\frac{1}{2}} \widetilde{M} \widetilde{D}^{-\frac{1}{2}}\).

After training, the anomaly score of each node \(x \in U\) is defined as the predicted probability of
belonging to the anomalous class:
\begin{equation}
S(x) = Z_{x2},
\end{equation}
where \(Z\) is the output of GCN, and \(Z_{x2}\) is the second element of the output probability vector
for node \(x\). A higher score indicates a greater likelihood of being an anomaly.

To determine the final set of outliers, an automatic threshold selection strategy based on the F1-score
is adopted. The nodes are first sorted in descending order of their anomaly scores, and candidate
thresholds are then evaluated at the top 1\% position, subsequently every 5\% up to 10\%, and every
10\% beyond 50\%. Since anomalies typically constitute only a small fraction of the dataset, the most
suspicious candidates are concentrated on by this search strategy, while computational efficiency is
maintained. The threshold that yields the highest \(F_{1}\) score on the labeled validation set is
finally selected for outlier identification.

\subsection{Illustrative Example}

To illustrate the proposed framework step by step, this subsection provides an example with the
information system \(IS = (U, A, V, f)\) in Table~\ref{tab:example_is}, where
\(U = \{x_{1}, x_{2}, x_{3}, x_{4}, x_{5}\}\), \(A = \{a, b, c\}\).

\begin{table}[!t]
\caption{The information system \(IS = (U, A, V, f)\)}
\label{tab:example_is}
\centering
\begin{tabular}{cccc}
\toprule
\(U\) & \(a\) & \(b\) & \(c\) \\
\midrule
\(x_{1}\) & 0.2 & 0.8 & X \\
\(x_{2}\) & 0.3 & 0.7 & X \\
\(x_{3}\) & 0.9 & 0.2 & Y \\
\(x_{4}\) & 0.1 & 0.9 & X \\
\(x_{5}\) & 0.9 & 0.3 & Y \\
\bottomrule
\end{tabular}
\end{table}

\textbf{(1) Compute the mean and standard deviation of numerical attributes:}
\(\mu_{a} = 0.48\), \(\sigma_{a} = 0.3487\); \(\mu_{b} = 0.58\), \(\sigma_{b} = 0.2786\);
\(\sigma_{c} = 0\).

\textbf{(2) Compute the attribute neighborhood entropy for each attribute.} Taking attribute \(a\) as
an example, the objects \(U\) are sorted by their values of \(a\):
\[
x_{4}(0.1) \leq x_{1}(0.2) \leq x_{2}(0.3) \leq x_{3}(0.9) \leq x_{5}(0.9).
\]
The gaps between consecutive values are:
\[
0.2-0.1 = 0.1,\; 0.3-0.2 = 0.1,\; 0.9-0.3 = 0.6,\; 0.9-0.9 = 0.
\]
With threshold \(\sigma_{a} = 0.3487\), the sorted sequence is split at the gap \(0.6 > 0.3487\),
yielding two connected components:
\[
C_{1} = \{x_{1}, x_{2}, x_{4}\},\quad C_{2} = \{x_{3}, x_{5}\}.
\]
By Definition~2,
\(NE(a) = -\frac{3}{5}\log_{2}\frac{3}{5} - \frac{2}{5}\log_{2}\frac{2}{5} \approx 0.9710\).

Similarly,
\(NE(b) = -\frac{2}{5}\log_{2}\frac{2}{5} - \frac{3}{5}\log_{2}\frac{3}{5} \approx 0.9710\),
and for categorical attribute \(c\),
\(NE(c) = -\frac{3}{5}\log_{2}\frac{3}{5} - \frac{2}{5}\log_{2}\frac{2}{5} \approx 0.9710\).

\textbf{(3) Compute the adaptive radii for each attribute by Definition~2:}
\begin{align*}
\varepsilon_{a} &= \frac{\sigma_{a}}{1 + \rho_{a}} = \frac{0.3487}{1 + 0.5818} \approx 0.2204, \\
\varepsilon_{b} &= \frac{\sigma_{b}}{1 + \rho_{b}} = \frac{0.2786}{1 + 0.5818} \approx 0.1761, \\
\varepsilon_{c} &= 0,
\end{align*}
in which
\[
\rho_{a} = 1 - \frac{NE(a)}{\log_{2}|U|} = 1 - \frac{0.9710}{2.3219} = 0.5818,\quad
\rho_{b} = 1 - \frac{NE(b)}{\log_{2}|U|} = 1 - \frac{0.9710}{2.3219} = 0.5818,\quad
\rho_{c} = 0.
\]

By Definition~1 and the adaptive radii above, the neighborhood knowledge is obtained and summarized in
Table~\ref{tab:neighborhood}.

\begin{table}[!t]
\caption{The neighborhood knowledge of \(IS\) in Table~1}
\label{tab:neighborhood}
\centering
\begin{tabular}{cccc}
\toprule
& \(N^{\varepsilon_{a}}(x)\) & \(N^{\varepsilon_{b}}(x)\) & \(N^{\varepsilon_{c}}(x)\) \\
\midrule
\(x_{1}\) & \(x_{1},x_{2},x_{4}\) & \(x_{1},x_{2},x_{4}\) & \(x_{1},x_{2},x_{4}\) \\
\(x_{2}\) & \(x_{1},x_{2},x_{4}\) & \(x_{1},x_{2}\) & \(x_{1},x_{2},x_{4}\) \\
\(x_{3}\) & \(x_{3},x_{5}\) & \(x_{3},x_{5}\) & \(x_{3},x_{5}\) \\
\(x_{4}\) & \(x_{1},x_{2},x_{4}\) & \(x_{1},x_{4}\) & \(x_{1},x_{2},x_{4}\) \\
\(x_{5}\) & \(x_{3},x_{5}\) & \(x_{3},x_{5}\) & \(x_{3},x_{5}\) \\
\bottomrule
\end{tabular}
\end{table}

\textbf{(4) Compute neighborhood information entropy for each object:}
\begin{align*}
NE(x_{1}) &= (-\tfrac{3}{5}\log_{2}\tfrac{3}{5}) + (-\tfrac{3}{5}\log_{2}\tfrac{3}{5}) + (-\tfrac{3}{5}\log_{2}\tfrac{3}{5}) \approx 1.3266,\\
NE(x_{2}) &= (-\tfrac{3}{5}\log_{2}\tfrac{3}{5}) + (-\tfrac{2}{5}\log_{2}\tfrac{2}{5}) + (-\tfrac{3}{5}\log_{2}\tfrac{3}{5}) \approx 1.4132,\\
NE(x_{3}) &= (-\tfrac{2}{5}\log_{2}\tfrac{2}{5}) + (-\tfrac{2}{5}\log_{2}\tfrac{2}{5}) + (-\tfrac{2}{5}\log_{2}\tfrac{2}{5}) \approx 1.5864,\\
NE(x_{4}) &= (-\tfrac{3}{5}\log_{2}\tfrac{3}{5}) + (-\tfrac{2}{5}\log_{2}\tfrac{2}{5}) + (-\tfrac{3}{5}\log_{2}\tfrac{3}{5}) \approx 1.4132,\\
NE(x_{5}) &= (-\tfrac{2}{5}\log_{2}\tfrac{2}{5}) + (-\tfrac{2}{5}\log_{2}\tfrac{2}{5}) + (-\tfrac{2}{5}\log_{2}\tfrac{2}{5}) \approx 1.5864.
\end{align*}

\textbf{(5)} By Definition~5, the neighborhood mutual information between all pairs of objects is
computed, and the results are shown in Table~\ref{tab:nmi}.

\begin{table}[!t]
\caption{Neighborhood mutual information}
\label{tab:nmi}
\centering
\begin{tabular}{cccccc}
\toprule
& \(x_{1}\) & \(x_{2}\) & \(x_{3}\) & \(x_{4}\) & \(x_{5}\) \\
\midrule
\(x_{1}\) & 1.3266 & 1.1792 & 0 & 1.1792 & 0 \\
\(x_{2}\) & 1.1792 & 1.4132 & 0 & 0.9488 & 0 \\
\(x_{3}\) & 0 & 0 & 1.5864 & 0 & 1.5864 \\
\(x_{4}\) & 1.1792 & 0.9488 & 0 & 1.4132 & 0 \\
\(x_{5}\) & 0 & 0 & 1.5864 & 0 & 1.5864 \\
\bottomrule
\end{tabular}
\end{table}

\textbf{(6)} Edge weights are computed by Eq.~(15). The resulting weight matrix is shown in
Table~\ref{tab:weights}, and the corresponding graph is illustrated in Fig.~\ref{fig:example_graph}.

\begin{table}[!t]
\caption{Normalized neighborhood mutual information (edge weights)}
\label{tab:weights}
\centering
\begin{tabular}{cccccc}
\toprule
& \(x_{1}\) & \(x_{2}\) & \(x_{3}\) & \(x_{4}\) & \(x_{5}\) \\
\midrule
\(x_{1}\) & 1 & 0.8612 & 0 & 0.8612 & 0 \\
\(x_{2}\) & 0.8612 & 1 & 0 & 0.6713 & 0 \\
\(x_{3}\) & 0 & 0 & 1 & 0 & 1 \\
\(x_{4}\) & 0.8612 & 0.6713 & 0 & 1 & 0 \\
\(x_{5}\) & 0 & 0 & 1 & 0 & 1 \\
\bottomrule
\end{tabular}
\end{table}

\begin{figure}[!t]
\centering
\includegraphics[width=0.7\columnwidth]{example_graph.pdf}
\caption{The constructed NMI graph.}
\label{fig:example_graph}
\end{figure}

% ===================================================================
\section{Experiments}
% ===================================================================
[To be completed: comprehensive experimental evaluation including data preprocessing,
comparison algorithms, parameter settings, and detailed analysis of results.]

\subsection{Data Preprocessing}
[To be completed.]

\subsection{Comparison Algorithms and Parameter Settings}
[To be completed.]

\subsection{Experimental Results and Analysis}
[To be completed: Precision, Recall, F1-score, AUC comparison tables; ROC/F1 curve analysis;
ablation study; statistical significance tests.]

\subsection{Parameter Sensitivity Analysis}
[To be completed.]

\subsection{Discussion}
[To be completed.]

% ===================================================================
\section{Conclusion}
% ===================================================================
[To be completed.]

% ===================================================================
\section*{Acknowledgments}
% ===================================================================
[To be completed.]

% ===================================================================
% References
% ===================================================================
\begin{thebibliography}{99}
\bibliographystyle{IEEEtran}

\bibitem{ref1}
V.~Chandola, A.~Banerjee, and V.~Kumar, ``Anomaly detection: A survey,''
\textit{ACM Comput. Surv.}, vol.~41, no.~3, pp.~1--58, 2009.

\bibitem{ref2}
C.~C.~Aggarwal, \textit{Outlier Analysis}, 2nd~ed. Springer, 2017.

\bibitem{ref3}
E.~W.~Ngai, Y.~Hu, Y.~H.~Wong, Y.~Chen, and X.~Sun, ``The application of data mining
techniques in financial fraud detection,'' \textit{Decis. Support Syst.}, vol.~50, no.~3,
pp.~559--569, 2011.

\bibitem{ref4}
A.~L.~Buczak and E.~Guven, ``A survey of data mining and machine learning methods for
cyber security intrusion detection,'' \textit{IEEE Commun. Surveys Tuts.}, vol.~18, no.~2,
pp.~1153--1176, 2016.

\bibitem{ref5}
P.~J.~Rousseeuw and A.~M.~Leroy, \textit{Robust Regression and Outlier Detection}.
John Wiley \& Sons, 1987.

\bibitem{ref6}
P.~Cunningham and S.~J.~Delany, ``$K$-nearest neighbour classifiers (with Python examples),''
\textit{arXiv preprint arXiv:2004.04523}, 2020.

\bibitem{ref7}
M.~M.~Breunig, H.-P.~Kriegel, R.~T.~Ng, and J.~Sander, ``LOF: Identifying density-based
local outliers,'' in \textit{Proc. ACM SIGMOD}, 2000, pp.~93--104.

\bibitem{ref8}
Z.~He, X.~Xu, and S.~Deng, ``Discovering cluster-based local outliers,''
\textit{Pattern Recognit. Lett.}, vol.~24, no.~9--10, pp.~1641--1650, 2003.

\bibitem{ref9}
G.~Pang, L.~Cao, L.~Chen, and H.~Liu, ``Unsupervised feature selection for outlier detection
by modelling hierarchical value-feature couplings,'' in \textit{Proc. ICDM}, 2016, pp.~410--419.

\bibitem{ref10}
J.~He, Q.~Xu, Y.~Jiang, Z.~Wang, and Q.~Huang, ``ADA-GAD: Anomaly-denoised autoencoders
for graph anomaly detection,'' in \textit{Proc. AAAI}, vol.~38, 2024, pp.~8481--8489.

\bibitem{ref11}
Z.~Pawlak, ``Rough sets,'' \textit{Int. J. Comput. Inf. Sci.}, vol.~11, no.~5, pp.~341--356, 1982.

\bibitem{ref12}
T.~Y.~Lin, ``Neighborhood systems and relational databases,'' in \textit{Proc. 16th ACM Annu.
Conf. Comput. Sci.}, 1988, p.~725.

\bibitem{ref13}
Q.~H.~Hu, D.~R.~Yu, J.~F.~Liu, and C.~X.~Wu, ``Neighborhood rough set based heterogeneous
feature subset selection,'' \textit{Inf. Sci.}, vol.~178, no.~18, pp.~3577--3594, 2008.

\bibitem{ref14}
Z.~Yuan, X.~Zhang, and S.~Feng, ``Hybrid data-driven outlier detection based on neighborhood
information entropy and its developmental measures,'' \textit{Expert Syst. Appl.}, vol.~112,
pp.~243--257, 2018.

\bibitem{ref15}
D.~R.~Wilson and T.~R.~Martinez, ``Improved heterogeneous distance functions,''
\textit{J. Artif. Intell. Res.}, vol.~6, no.~1, pp.~1--34, 1997.

\bibitem{ref16}
Z.~Yuan, H.~Chen, T.~Li, J.~Liu, and S.~Wang, ``Anomaly detection based on fuzzy neighborhood
rough sets,'' \textit{Inf. Sci.}, vol.~709, p.~122075, 2025.

\bibitem{ref17}
R.~Wille, ``Restructuring lattice theory: An approach based on hierarchies of concepts,''
in \textit{Ordered Sets}. Springer, 1982.

\bibitem{ref18}
B.~Ganter and R.~Wille, \textit{Formal Concept Analysis: Mathematical Foundations}. Springer, 2012.

\bibitem{ref19}
Q.~Hu, K.~Qin, L.~Zhang, and J.~Li, ``A novel outlier detection approach based on formal
concept analysis,'' \textit{Knowl.-Based Syst.}, vol.~268, p.~110486, 2023.

\bibitem{ref20}
J.~Li, K.~Qin, L.~Zhang, and Q.~Hu, ``Dual-aspect synergistic outlier detection with
structural deviation and attribute rarity,'' \textit{Pattern Recognit.}, vol.~180, p.~114084, 2026.

\bibitem{ref21}
T.~N.~Kipf and M.~Welling, ``Semi-supervised classification with graph convolutional networks,''
in \textit{Proc. ICLR}, 2017.

\bibitem{ref22}
A.~Goodge, B.~Hooi, S.~K.~Ng, and W.~S.~Ng, ``LUNAR: Unifying local outlier detection methods
via graph neural networks,'' in \textit{Proc. AAAI}, vol.~36, no.~6, 2022, pp.~6737--6745.

\bibitem{ref23}
K.~Ding, J.~Li, R.~Bhanushali, and H.~Liu, ``Deep anomaly detection on attributed networks,''
in \textit{Proc. SDM}, 2019, pp.~594--602.

\bibitem{ref24}
X.~Du, J.~Li, Z.~Chen, and Y.~Wang, ``Graph autoencoder-based unsupervised outlier detection,''
\textit{Inf. Sci.}, vol.~608, pp.~532--550, 2022.

\bibitem{ref25}
Q.~Li, Z.~Han, and X.-M.~Wu, ``Deeper insights into graph convolutional networks for
semi-supervised learning,'' in \textit{Proc. AAAI}, vol.~32, no.~1, 2018.

\bibitem{ref26}
K.~Liu, Y.~Dou, Y.~Zhao, X.~Ding, X.~Hu, R.~Zhang, K.~Ding, C.~Chen, H.~Peng, K.~Shu,
\textit{et al.}, ``Graph-based anomaly detection and deep learning: A comprehensive survey,''
\textit{ACM Comput. Surv.}, vol.~55, no.~8, pp.~1--37, 2022.

\bibitem{ref27}
F.~Jiang, Y.~Sui, and C.~Cao, ``An information entropy-based approach to outlier detection
in rough sets,'' \textit{Expert Syst. Appl.}, vol.~37, no.~9, pp.~6338--6344, 2010.

\bibitem{ref28}
G.~Pang, C.~Shen, L.~Cao, and A.~V.~D.~Hengel, ``Deep learning for anomaly detection: A review,''
\textit{ACM Comput. Surv.}, vol.~54, no.~2, pp.~1--38, 2021.

\bibitem{ref29}
G.~Duan, H.~Lv, H.~Wang, and G.~Feng, ``Application of a dynamic line graph neural network
for intrusion detection,'' \textit{IEEE Trans. Inf. Forensics Secur.}, vol.~18, pp.~699--714, 2023.

\bibitem{ref30}
H.~Yuan, H.~Yu, S.~Gui, and S.~Ji, ``Explainability in graph neural networks: A taxonomic survey,''
\textit{IEEE Trans. Pattern Anal. Mach. Intell.}, vol.~45, no.~5, pp.~6352--6372, 2022.

\bibitem{ref31}
S.~M.~Lundberg and S.-I.~Lee, ``A unified approach to interpreting model predictions,''
in \textit{Proc. NeurIPS}, vol.~30, 2017.

\bibitem{ref32}
W.~Hamilton, Z.~Ying, and J.~Leskovec, ``Inductive representation learning on large graphs,''
in \textit{Proc. NeurIPS}, vol.~30, 2017.

\bibitem{ref33}
M.~Feurer and F.~Hutter, ``Hyperparameter optimization,'' in \textit{Automated Machine Learning},
Springer, 2019, pp.~3--33.

\bibitem{ref34}
C.~Z.~Wang, Q.~H.~Hu, X.~Z.~Wang, D.~G.~Chen, and Y.~H.~Qian, ``Feature subset selection
based on fuzzy neighborhood rough sets,'' \textit{Knowl.-Based Syst.}, vol.~111, pp.~173--179, 2016.

\bibitem{ref35}
X.~Su, J.~Zhou, Y.~Chen, and W.~Pedrycz, ``Identifying outliers via local granular-ball density,''
\textit{IEEE Trans. Neural Netw. Learn. Syst.}, vol.~36, no.~10, pp.~18956--18967, 2025.

\bibitem{ref36}
J.~Yang, X.~Yang, J.~Zhou, and Y.~Yao, ``Three-way outlier detection based on shadowed
granular-balls,'' \textit{IEEE Trans. Fuzzy Syst.}, vol.~34, no.~1, pp.~101--113, 2026.

\bibitem{ref37}
A.~Roy, J.~Li, and S.~Pan, ``GAD-NR: Graph anomaly detection via neighborhood reconstruction,''
in \textit{Proc. WSDM}, 2024, pp.~576--585.

\bibitem{ref38}
N.~Rahimi and R.~Javadi, ``Application of graph autoencoder in outlier detection,''
\textit{Appl. Soft Comput.}, 2026.

\bibitem{ref39}
Z.~Qin, L.~Zhang, and K.~Li, ``Enhancing intrusion detection performance using GCN-LOF,''
\textit{Comput. Netw.}, 2025.

\bibitem{ref40}
Y.~Y.~Yao, ``Relational interpretations of neighborhood operators,''
\textit{Inf. Sci.}, vol.~111, no.~1--4, pp.~239--259, 1998.

\bibitem{ref41}
Y.~M.~Chen, D.~Q.~Miao, and R.~Z.~Wang, ``Neighborhood outlier detection,''
\textit{Expert Syst. Appl.}, vol.~37, no.~12, pp.~8745--8749, 2010.

\bibitem{ref42}
L.~A.~Zadeh, ``Fuzzy sets and information granularity,'' in \textit{Advances in Fuzzy Set Theory
and Applications}, 1979, pp.~3--18.

\bibitem{ref43}
D.~Dubois and H.~Prade, ``Rough fuzzy sets and fuzzy rough sets,''
\textit{Int. J. Gen. Syst.}, vol.~17, no.~2--3, pp.~191--209, 1990.

\bibitem{ref44}
Z.~Yuan, H.~Chen, T.~Li, B.~Sang, and S.~Wang, ``Outlier detection based on fuzzy rough granules
in mixed attribute data,'' \textit{IEEE Trans. Cybern.}, vol.~52, no.~8, pp.~8399--8412, 2021.

\bibitem{ref45}
Z.~Li, Y.~Zhao, J.~Han, Y.~Su, R.~Jiao, and X.~Wen, ``A survey on explainable anomaly detection,''
\textit{ACM Trans. Knowl. Discov. Data}, vol.~18, no.~1, pp.~1--54, 2024.

\bibitem{Mlyahilu2025}
N.~J.~Mlyahilu and E.~Novikova, ``Graph construction approaches for graph neural networks-based
anomaly detection in time series,'' preprint, 2025.

\bibitem{Febrinanto2025}
F.~G.~Febrinanto, K.~Moore, C.~Thapa, M.~Liu, V.~Saikrishna, J.~Ma, and F.~Xia,
``Entropy causal graphs for multivariate time series anomaly detection,''
\textit{ACM Trans. Intell. Syst. Technol.}, vol.~16, no.~6, Article~125, 2025.

\bibitem{Pang2024}
Asymptotic consistent graph structure learning for multivariate time-series anomaly detection, 2024.

\bibitem{He2024}
S.~He, G.~Li, K.~Xie, and P.~K.~Sharma, ``Fusion graph structure learning-based multivariate
time series anomaly detection with structured prior knowledge,''
\textit{IEEE Trans. Inf. Forensics Secur.}, vol.~19, pp.~8760--8772, 2024.

\end{thebibliography}

\end{document}
